\section{Lemezkezelés, particionálás és LUKS}

Részletes lépések a partíciók létrehozására, fájlrendszer készítésére és LUKS titkosításra.

\subsection{Particionálás és fájlrendszer}
\begin{lstlisting}[language=bash]
# Create GPT partition and ext4 filesystem (/dev/sdc1)
sudo parted /dev/sdc --script mklabel gpt mkpart primary ext4 1MiB 100%
sudo partprobe /dev/sdc
sudo mkfs.ext4 -L data /dev/sdc1
sudo mkdir -p /mnt/data
sudo mount /dev/sdc1 /mnt/data
\end{lstlisting}

\subsection{LUKS titkosítás}
\begin{lstlisting}[language=bash]
# Format and open LUKS volume
sudo cryptsetup luksFormat /dev/sdb1
sudo cryptsetup luksOpen /dev/sdb1 secret
sudo mkfs.xfs /dev/mapper/secret
sudo mount /dev/mapper/secret /mnt/secret
# Use a key file
sudo dd if=/dev/random of=/root/keyfile bs=16 count=1
sudo chmod 000 /root/keyfile
sudo cryptsetup luksAddKey /dev/sdb1 /root/keyfile
\end{lstlisting}

Megjegyzés: mindig készíts biztonsági mentést, és ellenőrizd a \texttt{blkid} és \texttt{lsblk -f} kimeneteket mielőtt \texttt{fstab}-ot szerkesztesz.
